\chapter{Testes e Análise de Resultados}

Este capítulo apresenta os testes realizados até ao momento, focados nas duas áreas já validadas de forma mais consistente: o comportamento do modelo físico impresso e a integração dos movimentos representados no protótipo.

\section{Testes do modelo físico}
Os primeiros testes confirmaram se o modelo impresso mantém a geometria pretendida e se permite a passagem adequada do broncoscópio, observando-se ainda deformações, zonas frágeis, arestas, dificuldade de limpeza e resistência ao manuseamento. No caso de materiais flexíveis, foi também avaliada a relação entre flexibilidade e estabilidade dimensional.

O teste inicial em PLA serviu como prova de conceito da geometria e da montagem, ainda que, por ser um material rígido, não representasse a solução final para realismo mecânico. Por isso, os ensaios posteriores com TPU, TPE ou outros filamentos flexíveis compararam facilidade de impressão, acabamento, deformação elástica e resistência à repetição dos movimentos.

Foi ainda realizada uma sessão preliminar de observação com um profissional de saúde, em que o broncoscópio foi introduzido no modelo impresso em 3D enquanto o sistema estava ativo (\cref{fig:teste-medico}). Esta sessão permitiu obter uma primeira avaliação qualitativa da disposição dos ramos, do espaço disponível para a progressão da sonda e da perceção de movimento durante o funcionamento do manipulador. As observações desta sessão ajudaram a identificar aspetos a melhorar na geometria do modelo e no comportamento do sistema, ainda que os dados quantitativos não tenham sido recolhidos de forma sistemática.

\begin{figure}[htbp]
  \centering
  \includegraphics[width=0.82\textwidth]{fotos/teste_medico_broncoscopio.jpg}
  \caption{Sessão de observação preliminar com profissional de saúde: broncoscópio introduzido no modelo de árvore brônquica impresso em 3D enquanto o sistema de controlo estava ativo.}
  \label{fig:teste-medico}
\end{figure}

\section{Testes de integração dos movimentos}
A integração foi feita por etapas. Primeiro executou-se apenas a respiração, verificando periodicidade, amplitude e suavidade. Depois foi adicionado o batimento cardíaco, que deve aparecer como pequena oscilação sobreposta. Por fim, testou-se a tosse, que altera temporariamente o comportamento da respiração e ativa uma vibração ou perturbação curta.

O modo combinado é o mais importante para avaliar o realismo, pois aproxima o comportamento de um ambiente broncoscópico dinâmico. Neste modo, foram observados conflitos entre movimentos, saturação dos servos, perda de suavidade e eventuais posições impossíveis da cinemática inversa.
